\section{Physiological constraints on receptor-mediated ligand uptake}\label{sec:physiological}
Having validated the analytical model against RW simulation, we now use the theory to examine the behavior of the system under physiological conditions.
Specifically, we substitute representative physiological parameter values into the analytical expressions to assess how ligand uptake depends on receptor abundance and blood concentration $C_{0}$.

\paragraph{Flux efficiency.}
We focus on the normalized uptake efficiency $\mu = J/\Jmax$ as a function of receptor surface coverage for different ligand concentrations.
As shown in Fig.~\ref{fig:mu_coverage}, the uptake efficiency at physiological receptor numbers (on the order of $10^{5}$--$10^{6}$ receptors) is significantly below 1.
For example, at $N \sim 10^{5}$ receptors, $\mu \approx 5 \times 10^{-3}$, indicating that ligand uptake is far from diffusion-limited.

This behavior reflects the fact that insulin receptor kinetics are intrinsically slow, with characteristic rates on the order of $10^{-3}\,\mathrm{s}^{-1}$.
As a result, uptake is limited primarily by receptor kinetics rather than by diffusive transport to the cell surface.
In this regime, the steady-state concentration profile remains nearly flat, $C(r) \simeq C_{0}$, consistent with the limit $\mu \ll 1$ in Eq.~\eqref{eq:mu}.

Importantly, at low physiological ligand concentrations ($C_{0} \sim 1$--$100$\,pM), the uptake efficiency is higher and largely concentration-independent, reflecting unsaturated receptor kinetics.
At higher concentrations, receptor saturation becomes significant, reducing uptake efficiency as an increasing fraction of receptors remains occupied.

\begin{figure}[H]
    \centering
    \fbox{\parbox[c][7cm][c]{0.75\textwidth}{\centering \textit{Placeholder for Figure 5: $\mu = J/\Jmax$ vs.\ receptor surface coverage.}}}
    \caption{\textbf{Normalized ligand uptake flux.} The normalized flux $J/\Jmax$, computed from Eq.~\eqref{eq:mu_quadratic}, is shown as a function of receptor surface coverage on the outer cylinder for different ligand concentrations $C_{0}$. Physiological parameter values are used throughout.}
    \label{fig:mu_coverage}
\end{figure}

\paragraph{Receptor occupancy.}
Beyond the total ligand uptake flux, a key quantity of biological relevance is the receptor occupancy, since intracellular signaling is initiated by ligand-bound (and subsequently phosphorylated) receptors at the cell membrane.
In particular, downstream signaling pathways that regulate GLUT1 expression are driven by the fraction of occupied receptors rather than by the absolute ligand flux.

To examine this behavior, we compute the steady-state fraction of bound receptors as a function of the ligand concentration in the blood, $C_{0}$, for different total numbers of receptors (Fig.~\ref{fig:occupancy}).
One might expect that increasing the number of receptors would reduce the occupancy fraction, due to enhanced ligand consumption and a corresponding depletion of ligand near the cell surface.

Surprisingly, we find that the occupancy curves are nearly identical for different receptor abundances.
This behavior is a direct consequence of the fact that the uptake efficiency $\mu$ is much smaller than unity under physiological conditions.
As discussed above, when $\mu \ll 1$, ligand depletion at the cell surface is negligible and the local concentration at the receptors satisfies $C(R) \simeq C_{0}$.
Consequently, receptor occupancy is governed primarily by intrinsic receptor kinetics and depends only weakly on the total number of receptors.
For physiological insulin concentrations (10--500\,pM), the predicted receptor occupancy ranges from $\sim 0.3\%$ to 15\%.
Under basal conditions ($C_{0} \approx 50\,\mathrm{pM}$), occupancy is only $\sim 2\%$, indicating that receptors operate far from saturation.

\begin{figure}[H]
    \centering
    \fbox{\parbox[c][7cm][c]{0.75\textwidth}{\centering \textit{Placeholder for Figure 6: Bound receptor percentage vs $C_{0}$.}}}
    \caption{\textbf{Receptor occupancy under physiological conditions.} The steady-state fraction of bound receptors is computed from Eq.~\eqref{eq:mm_relations} as a function of the blood ligand concentration $C_{0}$ for different total receptor numbers. Because ligand depletion at the cell surface is negligible ($\mu \ll 1$), the occupancy curves collapse onto the Michaelis--Menten prediction and are largely independent of receptor abundance.}
    \label{fig:occupancy}
\end{figure}
